%%% FIELD transfer-statement
Suppose \(q(v)>0\) for every \(v\in\Omega_V\).  Every recursive SCM \(M\) for which there is a \(\rho\in[\delta,1]\) satisfying \(c_v(M)=\rho q(v)\) for all \(v\in\Omega_V\) and \(d_M\geq\delta\) obeys the graph-free bounds
\[
L_\delta\leq\theta(M)\leq U_\delta.
\]
If \(\mathcal P\) is the complete all-circle COPAG, every observed variable has at least two states, and \(|\Omega_B|\geq2\), there is a one-source family \(\{M_t:0\leq t\leq1\}\subseteq\mathfrak M_\delta(\mathcal P,q)\) such that
\[
\theta(M_t)=(1-t)L_\delta+tU_\delta.
\]
Each endpoint witness has at most \(m+1\) positive-mass response types.  Consequently, for any family \(\mathscr C\) of recursive SCMs all satisfying the graph-free premises above and containing every \(M_t\),
\[
\{\theta(M):M\in\mathscr C\}=[L_\delta,U_\delta].
\]
In particular, the same equality holds for every permitted-graph family satisfying those premises and containing the complete-COPAG witnesses.
%%% FIELD transfer-proof
Put
\[
I_b=\{v\in\Omega_V:B(v)=b\},\qquad
R_- =\sum_{v\in I_b}q(v)\{h(A(v))-h_-\},\qquad
R_+ =\sum_{v\in I_b}q(v)\{h_+-h(A(v))\}.
\]
Finiteness of \(\Omega_A\) implies that the extrema defining \(h_-\) and \(h_+\) are attained.  Their definitions give \(R_-\geq0\) and \(R_+\geq0\).  Also
\[
U_\delta-L_\delta
=(h_+-h_-)\left\{1-\delta\sum_{v\in I_b}q(v)\right\}\geq0,
\]
because \(0<\delta\leq1\) and \(\sum_{v\in I_b}q(v)\leq1\).  Thus the claimed interval is nonempty.

First fix an arbitrary recursive SCM \(M\) satisfying the three graph-free premises in the statement.  The margin \(d_M\geq\delta>0\) makes the division in \(\theta(M)=n_M/d_M\) legitimate.  For every \(v\in I_b\), recursive consistency gives
\[
\{V=v,S=1\}\subseteq\{A^b=A(v),S^b=1\}. \tag{1}
\]
To verify (1), order the endogenous variables topologically.  On \(\{V=v,S=1\}\), replacing the structural equations for \(B\) by the already realized value \(b\) does not change \(B\).  Induction along the remaining structural equations leaves every later endogenous variable, including \(A\) and \(S\), unchanged.  The factual events on the left side of (1), indexed by \(v\in I_b\), are disjoint.  Since \(h(A^b)-h_-\geq0\), consistency (1), selected-law fit, and \(\rho\geq\delta\) imply
\[
\begin{aligned}
n_M-h_-d_M
&=\mathbb E_M\!\left[\{h(A^b)-h_-\}\mathbf 1\{S^b=1\}\right]\\
&\geq\sum_{v\in I_b}\{h(A(v))-h_-\}\Pr_M(V=v,S=1)\\
&=\rho R_-\\
&\geq\delta R_-.
\end{aligned} \tag{2}
\]
Because \(d_M\) is a probability and is positive, \(0<d_M\leq1\).  Dividing (2) by \(d_M\) therefore yields
\[
\theta(M)-h_-
=\frac{n_M-h_-d_M}{d_M}
\geq\frac{\delta R_-}{d_M}
\geq\delta R_-,
\]
which is \(\theta(M)\geq L_\delta\) by \(\text{def:saturated-envelope}\).  Applying the same argument to the nonnegative outcome map \(h_+-h\) gives
\[
\begin{aligned}
h_+d_M-n_M
&=\mathbb E_M\!\left[\{h_+-h(A^b)\}\mathbf 1\{S^b=1\}\right]\\
&\geq\sum_{v\in I_b}\{h_+-h(A(v))\}\Pr_M(V=v,S=1)\\
&=\rho R_+\\
&\geq\delta R_+.
\end{aligned} \tag{3}
\]
Division by \(d_M\leq1\) gives
\[
h_+-\theta(M)
=\frac{h_+d_M-n_M}{d_M}
\geq\frac{\delta R_+}{d_M}
\geq\delta R_+,
\]
or \(\theta(M)\leq U_\delta\).  This proves the outer bound without using any condition on \(G_S(M)\).

We now construct witnesses in the exact complete-COPAG subclass.  Choose \(a_-,a_+\in\Omega_A\) such that \(h(a_-)=h_-\) and \(h(a_+)=h_+\).  The hypothesis \(|\Omega_B|\geq2\) supplies a value \(b^\dagger\in\Omega_B\setminus\{b\}\).  Let one exogenous source \(U\) have a state \(u_v\) for every \(v\in\Omega_V\), with
\[
\Pr(U=u_v)=\delta q(v).
\]
When \(\delta<1\), add a state \(u_0\) of mass \(1-\delta\); when \(\delta=1\), omit that state.  Since \(q(v)>0\), all displayed \(u_v\)-states have positive mass, and the masses sum to one.  Hence each endpoint source has at most \(m+1\) positive-mass states.

Choose a recursive order in which the vertices of \(B\) precede those of \(V\setminus B\), which in turn precede \(S\).  Give the vector \(B\) its natural value \(B(v)\) at \(u_v\), and give it the natural value \(b^\dagger\) at \(u_0\), if that state exists.  For each \(X\in V\setminus B\), take a response table \(f_X(B,U)\) whose factual entries satisfy
\[
f_X(B(v),u_v)=X(v). \tag{4}
\]
For the selection table impose
\[
f_S(B(v),u_v)=1,qquad
f_S(b,u)=1\ \text{for every positive-mass state }u, \tag{5}
\]
and, if \(u_0\) exists, impose \(f_S(b^\dagger,u_0)=0\).  In the lower endpoint model, complete the intervention entries of the coordinates of \(A\) so that
\[
A(b,u_v)=
\begin{cases}
A(v),&B(v)=b,\\
a_-,&B(v)\neq b,
\end{cases}
\qquad
A(b,u_0)=a_-\quad(\delta<1). \tag{6-}
\]
For the upper endpoint model use
\[
A(b,u_v)=
\begin{cases}
A(v),&B(v)=b,\\
a_+,&B(v)\neq b,
\end{cases}
\qquad
A(b,u_0)=a_+\quad(\delta<1). \tag{6+}
\]
There is no conflict among (4)--(6): if \(B(v)=b\), the intervention entry is exactly the factual entry; every changed \(u_v\)-entry has \(B(v)\neq b\); and \(b^\dagger\neq b\).  Fill all remaining response-table entries arbitrarily.  The resulting finite tables define recursive one-source SCMs.

Under the natural regime, precisely the \(u_v\)-types are selected.  Therefore
\[
c_v(M)=\Pr_M(V=v,S=1)=\delta q(v)\quad(v\in\Omega_V),
\qquad \rho=\delta. \tag{7}
\]
By (5), every positive-mass type is selected after \(\operatorname{do}(B=b)\), and hence
\[
d_M=1\geq\delta. \tag{8}
\]
Equations (6-) and (8) give, for the lower model,
\[
\begin{aligned}
\theta(M)
=n_M
&=\delta\sum_{v\in I_b}q(v)h(A(v))
+\left\{1-\delta\sum_{v\in I_b}q(v)\right\}h_-\\
&=h_-+\delta\sum_{v\in I_b}q(v)\{h(A(v))-h_-\}\\
&=L_\delta.
\end{aligned} \tag{9-}
\]
Likewise, (6+) and (8) give
\[
\begin{aligned}
\theta(M)
=n_M
&=\delta\sum_{v\in I_b}q(v)h(A(v))
+\left\{1-\delta\sum_{v\in I_b}q(v)\right\}h_+\\
&=h_+-\delta\sum_{v\in I_b}q(v)\{h_+-h(A(v))\}\\
&=U_\delta.
\end{aligned} \tag{9+}
\]

It remains to prove exact graphical legality.  Fix \(X\in V\).  If \(X\in B\), the assumption that \(X\) has at least two states and positivity of every joint cell provide positive-mass source states at which the natural response of \(X\) differs.  If \(X\notin B\), choose \(v,v'\in\Omega_V\) with \(B(v)=B(v')\), \(X(v)\neq X(v')\), and all other coordinates fixed.  Such cells exist because the joint state space is a Cartesian product and \(X\) has at least two states; positivity of \(q\) makes both corresponding source states positive-mass.  Equation (4) yields
\[
f_X(B(v),u_v)=X(v)\neq X(v')=f_X(B(v),u_{v'}).
\]
Thus \(U\) is an actual functional-essential parent of every observed vertex, under the actual-parent convention defining \(G_S(M)\), with all other endogenous inputs held fixed.  Consequently, for every distinct \(X,Y\in V\), the latent-fork path \(X\leftarrow U\to Y\) is present and remains open after conditioning on \(S\) and on any subset of \(V\setminus\{X,Y\}\): its only nonendpoint, \(U\), is an unconditioned noncollider.  The selected maximal projection therefore makes every pair of observed vertices adjacent.  Its selected MAG has complete skeleton.  A complete-skeleton selected MAG has no unshielded collider and no discriminating path with nonadjacent endpoints, so it belongs to the Markov-equivalence class represented by the complete all-circle COPAG.  Hence
\[
G_S(M)\in\mathcal G(\mathcal P). \tag{10}
\]
Equations (7), (8), and (10), together with \(\rho=\delta\in[\delta,1]\), prove that both endpoint SCMs belong to \(\mathfrak M_\delta(\mathcal P,q)\).

For \(0<t<1\), take disjointly labelled copies of the lower and upper source states, multiply their masses by \(1-t\) and \(t\), and use the corresponding response tables on the two copies.  At \(t=0\) and \(t=1\), use the respective endpoint SCM.  This still gives a recursive SCM with one exogenous source.  Equations (7) and (8) are preserved by mixing, and the positive-mass essential-source argument proving (10) remains valid.  Thus \(M_t\in\mathfrak M_\delta(\mathcal P,q)\) for every \(t\in[0,1]\).  Since its intervention denominator is one and expectation is linear in the source law,
\[
\theta(M_t)=(1-t)L_\delta+tU_\delta. \tag{11}
\]

Finally let \(\mathscr C\) satisfy the hypotheses in the statement.  The graph-free bound already proved gives
\[
\{\theta(M):M\in\mathscr C\}\subseteq[L_\delta,U_\delta].
\]
Because \(\mathscr C\) contains every \(M_t\), equation (11) gives the reverse inclusion:
\[
[L_\delta,U_\delta]
=\{\theta(M_t):0\leq t\leq1\}
\subseteq\{\theta(M):M\in\mathscr C\}.
\]
This proves the family-sandwich assertion, including its permitted-graph specialization.
%%% FIELD target-proof
Apply \(\text{thm:saturated-envelope-transfer}\) to the complete all-circle COPAG in the theorem.  Its graph-free implication applies to every \(M\in\mathfrak M_\delta(\mathcal P,q)\), because \(\text{def:admissible-class}\) supplies a common \(\rho\in[\delta,1]\) with \(c_v(M)=\rho q(v)\) for all \(v\), as well as \(d_M\geq\delta\).  Therefore
\[
\Theta_\delta(\mathcal P,q)\subseteq[L_\delta,U_\delta]. \tag{1}
\]
The same transfer theorem constructs, using one exogenous source, models \(M_t\in\mathfrak M_\delta(\mathcal P,q)\) satisfying
\[
\theta(M_t)=(1-t)L_\delta+tU_\delta,
\qquad 0\leq t\leq1. \tag{2}
\]
Equation (2) proves the reverse inclusion in (1).  It also states that the endpoint models \(M_0\) and \(M_1\) have at most \(m+1\) positive-mass response types.  Hence
\[
\Theta_\delta(\mathcal P,q)=[L_\delta,U_\delta]
\]
with both endpoints attained as claimed.

If \(h\) is constant, then \(h_-=h_+\) and, for every compatible model,
\[
n_M=\mathbb E_M\!\left[h_-\mathbf 1\{S^b=1\}\right]=h_-d_M.
\]
The margin \(d_M\geq\delta>0\) permits division, giving \(\theta(M)=h_-\).  The witness construction in \(\text{thm:saturated-envelope-transfer}\) proves nonemptiness under the nonvacuous hypotheses, so the fiber is \(\{h_-\}\).

Finally, if the intervention is vacuous, its submodel is the factual model.  Selected-law fit and normalization of \(q\) give
\[
d_M=\Pr_M(S=1)=\sum_{v\in\Omega_V}c_v(M)=\rho
\]
and
\[
n_M=\sum_{v\in\Omega_V}h(A(v))c_v(M)
=\rho\sum_{v\in\Omega_V}q(v)h(A(v)).
\]
The population margin \(\rho\geq\delta>0\) makes division legitimate, so every compatible SCM has
\[
\theta(M)=\sum_{v\in\Omega_V}q(v)h(A(v)). \tag{3}
\]
For nonemptiness in the complete-COPAG class, let one source have states \(u_v\) of masses \(\delta q(v)\), together, when \(\delta<1\), with one state \(u_0\) of mass \(1-\delta\).  Put \(V(u_v)=v\), \(S(u_v)=1\), and \(S(u_0)=0\), choosing \(V(u_0)\) arbitrarily.  Then \(c_v=\delta q(v)\), \(\rho=d_M=\delta\), and strict positivity of \(q\), together with at least two states per observed variable, makes the common source an actual functional-essential parent of every observed vertex.  The latent-fork argument in \(\text{thm:saturated-envelope-transfer}\) therefore gives a complete selected MAG and exact membership in the complete all-circle COPAG class.  Thus a compatible model exists, and (3) proves the asserted singleton fiber.
